\chapter{Source Code Excerpts}
\label{app:code}

This appendix contains key source code excerpts from the implementation.

\section{Token Bucket Implementation}

\begin{lstlisting}[language=Kotlin, caption={TokenBucket.kt}]
class TokenBucket(
    @Volatile var rateBps: Long,
    @Volatile var burstBytes: Long
) {
    private var tokens: Double = burstBytes.toDouble()
    private var lastRefillTime: Long = System.nanoTime()

    @Synchronized
    fun consume(bytes: Int): Boolean {
        refill()
        return if (tokens >= bytes) {
            tokens -= bytes
            true
        } else {
            false
        }
    }

    @Synchronized
    fun setRate(newRateBps: Long) {
        rateBps = newRateBps
        burstBytes = newRateBps * 2
        tokens = minOf(tokens, burstBytes.toDouble())
    }

    private fun refill() {
        val now = System.nanoTime()
        val elapsed = (now - lastRefillTime) / 1_000_000_000.0
        tokens = minOf(burstBytes.toDouble(), tokens + rateBps * elapsed)
        lastRefillTime = now
    }
}
\end{lstlisting}

\section{DPI Classifier Implementation}

\begin{lstlisting}[language=Kotlin, caption={DpiClassifier.kt}]
class DpiClassifier {
    private val flowCache = ConcurrentHashMap<FlowKey, TrafficCategory>()

    fun classify(packet: RawPacket): TrafficCategory {
        val key = FlowKey(packet.srcIp, packet.dstIp, 
                         packet.srcPort, packet.dstPort, packet.protocol)
        return flowCache.getOrPut(key) { 
            matchRules(packet.dstPort, packet.protocol) 
        }
    }

    private fun matchRules(dstPort: Int, protocol: Protocol): TrafficCategory {
        return when {
            // VoIP
            protocol == Protocol.UDP && dstPort in listOf(5060, 5061) 
                -> TrafficCategory.VOIP
            protocol == Protocol.UDP && dstPort in 16384..32767 
                -> TrafficCategory.VOIP
            
            // Video Conferencing
            protocol == Protocol.UDP && dstPort == 3478 
                -> TrafficCategory.VIDEO_CONFERENCING
            protocol == Protocol.UDP && dstPort in 19302..19309 
                -> TrafficCategory.VIDEO_CONFERENCING
            
            // Online Gaming
            protocol == Protocol.UDP && dstPort in 27015..27030 
                -> TrafficCategory.ONLINE_GAMING
            protocol == Protocol.UDP && dstPort == 3074 
                -> TrafficCategory.ONLINE_GAMING
            
            // File Transfer
            protocol == Protocol.TCP && dstPort in listOf(20, 21, 22, 445, 2049) 
                -> TrafficCategory.FILE_TRANSFER
            
            // Web Browsing
            protocol == Protocol.TCP && dstPort in listOf(80, 443) 
                -> TrafficCategory.WEB_BROWSING
            
            else -> TrafficCategory.UNKNOWN
        }
    }
}
\end{lstlisting}

\section{Packet Parser Implementation}

\begin{lstlisting}[language=Kotlin, caption={RawPacket.kt - IPv4 Parsing}]
companion object {
    fun parse(buffer: ByteArray, length: Int): RawPacket? {
        if (length < 20) return null
        val buf = ByteBuffer.wrap(buffer, 0, length)
        val versionIhl = buf.get(0).toInt() and 0xFF
        val version = versionIhl shr 4

        return when (version) {
            4 -> parseIpv4(buf, buffer, length)
            6 -> parseIpv6(buf, buffer, length)
            else -> null
        }
    }

    private fun parseIpv4(buf: ByteBuffer, raw: ByteArray, length: Int): RawPacket? {
        val ihl = (buf.get(0).toInt() and 0x0F) * 4
        if (length < ihl + 4) return null

        val protocolNum = buf.get(9).toInt() and 0xFF
        val protocol = Protocol.fromNumber(protocolNum)

        val srcIp = buildString {
            for (i in 12..15) {
                if (i > 12) append('.')
                append(buf.get(i).toInt() and 0xFF)
            }
        }
        val dstIp = buildString {
            for (i in 16..19) {
                if (i > 16) append('.')
                append(buf.get(i).toInt() and 0xFF)
            }
        }

        val (srcPort, dstPort) = extractPorts(buf, ihl, protocol)

        return RawPacket(srcIp, dstIp, srcPort, dstPort, protocol, 
                        length, raw.copyOf(length))
    }

    private fun extractPorts(buf: ByteBuffer, headerEnd: Int, 
                            protocol: Protocol): Pair<Int, Int> {
        if (protocol == Protocol.OTHER) return Pair(0, 0)
        if (buf.capacity() < headerEnd + 4) return Pair(0, 0)
        val src = ((buf.get(headerEnd).toInt() and 0xFF) shl 8) or 
                  (buf.get(headerEnd + 1).toInt() and 0xFF)
        val dst = ((buf.get(headerEnd + 2).toInt() and 0xFF) shl 8) or 
                  (buf.get(headerEnd + 3).toInt() and 0xFF)
        return Pair(src, dst)
    }
}
\end{lstlisting}

\section{Bandwidth Scheduler Implementation}

\begin{lstlisting}[language=Kotlin, caption={BandwidthScheduler.kt - Rebalancing}]
fun rebalanceWithDevices(devices: List<ConnectedDevice>) {
    val highCount = devices.count { it.priorityClass == TrafficClass.HIGH }
    val medCount  = devices.count { it.priorityClass == TrafficClass.MEDIUM }
    val lowCount  = devices.count { it.priorityClass == TrafficClass.LOW }

    val totalWeight = highCount * 4 + medCount * 2 + lowCount * 1
    if (totalWeight == 0) return

    devices.forEach { device ->
        if (manualCaps.containsKey(device.ipAddress)) return@forEach
        
        val weight = when (device.priorityClass) {
            TrafficClass.HIGH   -> 4
            TrafficClass.MEDIUM -> 2
            TrafficClass.LOW    -> 1
        }
        val rate = (uplinkBps * weight) / totalWeight
        val burst = when (device.priorityClass) {
            TrafficClass.HIGH   -> rate * 2
            TrafficClass.MEDIUM -> rate
            TrafficClass.LOW    -> rate / 2
        }
        buckets.getOrPut(device.ipAddress) { TokenBucket(rate, burst) }
            .setRate(rate)
    }
}
\end{lstlisting}

\textit{(Complete source code available at: https://github.com/[repository-url])}
